\documentclass[10pt,a4paper]{amsart}
\usepackage[margin=19mm]{geometry}
\usepackage{amsmath,amssymb,mathtools}
\usepackage[scr=rsfs,cal=euler]{mathalfa}
\usepackage{xfrac}
\usepackage[unicode,hidelinks,bookmarks=false]{hyperref}
\newcommand{\ie}{i.e.\ }
\newcommand{\cf}{cf.\ }
\newcommand{\resp}{respectively\ }
\newcommand{\Subsection}{\subsection}
\providecommand{\nobreakdash}{\nobreak\hbox{-}}
\setlength{\emergencystretch}{2em}
\multlinegap0pt
\usepackage{amssymb}
\usepackage{mathtools}
\usepackage[normalem]{ulem}
\usepackage{xcolor}
\usepackage{tikz}
\newtheorem{lem}{Lemma}
\newcommand{\lan}{\langle }
\newcommand{\ran}{\rangle}
\title{New and old Saito-Kurokawa lifts classically via $L^2$ norms and bounds on their sup-norms: level aspect}
\author{Pramath Anamby, Soumya Das}
\date{Source: arXiv:2403.17401v3 (CC BY 4.0)}
\begin{document}
\maketitle

\noindent\emph{Source and licence.} New and old Saito-Kurokawa lifts classically via $L^2$ norms and bounds on their sup-norms: level aspect. Version arXiv:2403.17401v3 (CC BY 4.0), distributed by arXiv under the Creative Commons Attribution 4.0 International licence, http://creativecommons.org/licenses/by/4.0/, as recorded in the arXiv licence metadata for this exact version and shown on the abstract page https://arxiv.org/abs/2403.17401v3. Full licence notice: \texttt{licenses/arxiv-2403.17401-CC-BY-4.0.txt}.\par\smallskip

\noindent\textit{Editorial scope.} Counted unit: the article's lem lem:phialphabeta (source0.tex lines 2804--2806). The article's complete original proof of it is delivered with it (source0.tex lines 2857--2890, one \texttt{proof} environment of 34 lines, which the article places away from the statement). No mathematics is written, repaired or re-derived: the statement and the proof are the article's own lines, in the article's own notation, and no retained line is substituted or re-flowed.  Retained uncounted as the source-local context the proof needs: lines 2814--2816; lines 2849--2853.  Nothing was omitted from any retained span, so the package carries no omission record.  No retained line contains a citation, so the wrapper carries no bibliography and no bibliography entry.  No retained line contains a figure environment, an image-inclusion command or drawing code, so no image is delivered and the package declares no asset dependency.  Editorial formatting only, in addition: an AMS \texttt{amsart} preamble that restates the article's own declarations and macros used by the retained lines, namely the environments \texttt{lem} on separate counters, together with the standard packages those lines need.  The retained environments print in this document as Lemma 1 in order of appearance, whereas the article prints the same items with the numbering of its own full document.  The complete native source of the article as distributed in the arXiv e-print for this exact version is bundled unchanged as \texttt{sources/156649/source0.tex}.
\begin{align} \label{ujp-up2}
\lan\phi|U_J(p), \phi\ran_p = \frac{[\Gamma_0(p):\Gamma_0(4p)]}{2^{2k-1/2}}\lan h|U(p^2), h\ran_p.    
\end{align}

\begin{equation} \label{up2}
    \begin{split}
        \lan h|U(p^2) , g\ran_{4p^2}=  p^{k-5/2}\cdot p^2 \lan h|B(p^2), g\ran_{4p^2}= \frac{p\lambda_h(p)}{(p+1)}\lan h, g\ran_{4p^2}.
    \end{split}
\end{equation}

\begin{lem}\label{lem:phialphabeta}
   With the above notation, $\lan \phi_\alpha, \phi_\beta\ran_p \neq 0$.
\end{lem}

\begin{proof}[Proof of Lemma \ref{lem:phialphabeta}]
We first note that
\begin{align}\label{phialphabetaIP}
    \lan \phi_\alpha, \phi_\beta\ran_p&=\lan \phi,\phi\ran_p - p^{1-k}\alpha\lan \phi|w_p, \phi\ran_p-p^{1-k}\overline{\beta}\lan \phi, \phi|w_p\ran_p+ p^{2-2k}\alpha\overline{\beta}\lan \phi, \phi\ran_p.
\end{align}
Now we find an expression for $\lan \phi|w_p, \phi\ran_p$. Since $w_p= p^{2-k}(T_J(p)-U_J(p))$, and $\phi$ is an eigenfunction of $T_J(p)$,  
\begin{align}
\lan \phi|w_p, \phi\ran_p= p^{2-k}(\lambda_p\lan \phi, \phi\ran_p -\lan \phi|U_J(p), \phi\ran_p).
\end{align}
%
Now, combining \eqref{ujp-up2} and \eqref{up2}, we get,
\begin{align}
    \lan \phi|U_J(p), \phi\ran_p= \frac{p}{p+1}\lambda_p \lan\phi, \phi\ran_p.
\end{align}
As a consequence, we get
\begin{align} \label{phi-wp}
    \lan \phi|w_p, \phi\ran_p= \frac{p^{2-k}}{p+1}\lambda_p \lan \phi, \phi\ran_p.
\end{align}
This expression also implies that $\lan \phi|w_p, \phi\ran_p= \lan \phi, \phi|w_p\ran_p$, as they are real valued.


So if the inner product $ \lan \phi_\alpha, \phi_\beta\ran_p=0$ in \eqref{phialphabetaIP}, since $\overline{\beta}=\alpha$, $\alpha$ must satisfy the polynomial
\begin{equation}
    p^{2k-2}- 2\frac{p\lambda_p}{p+1} X+ X^2,
\end{equation}
since $\lan \phi|w_p, \phi\ran_p= \lan \phi, \phi|w_p\ran_p$ are given by \eqref{phi-wp}.

But $\alpha$ also satisfies the polynomial $X^2-\lambda_p X+ p^{2k-3}$. Thus, we get that
\begin{align}
    &\alpha\lambda_p\left(\frac{2p}{p+1}-1\right)+p^{2k-3}(1-p)=0\implies \alpha^2+p^{2k-3}= p^{2k-2}+p^{2k-3}.\\
    %&\alpha\lambda_p= p^{2k-2}+p^{2k-3}\\
\end{align}
That is, $|\alpha|^2= p^{2k-2}$, which is not possible since $|\alpha|^2=\alpha\beta=p^{2k-3}$ and this gives us that $\lan \phi_\alpha, \phi_\beta\ran_p\neq 0$, as desired.
\end{proof}

\end{document}
